\documentclass[../main.tex]{subfiles}

\begin{document}
	\chapter{Lista de funciones desarrolladas}
	\label{sec:ListaFuncionesDerarrolladas}
	
	Se listan a continuación las funciones de cada fichero JavaScript
	del proyecto, en el orden en que los carga \texttt{index.html},
	agrupadas por etapa: el puente de datos, la autoría (\texttt{Base*}),
	y el juego (\texttt{Juego*}). Para cada una se detalla su nombre y una
	s
	\textbf{ESTA LISTA ES PROVISIONAL}. Comprende unas 200 funciones desarrolladas hasta el fin de la fase 4 inclusive, se irá ampliando.
	
	\subsection*{comunListas.js — el puente de datos}
	
	\begin{itemize}
		\item \texttt{cargar}: guarda en el puente el setup y la lista de
		segmentos del GWI (filtra los fantasmas tipo 99).
		\item \texttt{limpiar}: vacía todo el puente (listas, fondo,
		configuraciones).
		\item \texttt{segmentos} / \texttt{sprites} / \texttt{fondo}:
		lectores de las referencias vivas (el juego lee; nunca escribe).
		\item \texttt{fijarFondo}: fija la imagen de fondo del cartucho.
		\item \texttt{fijarReglas} / \texttt{reglas}: escribe la autoría /
		lee el motor la regla de combinación de pulsaciones del setup.
		\item \texttt{fijarFin} / \texttt{finDeCartucho}: fija / lee el
		mecanismo de fin de partida del cartucho.
		\item \texttt{fijarSecuencia} / \texttt{secuenciaDeCartucho}: fija /
		lee la Secuencia (animación de fondo).
		\item \texttt{fijarPausas} / \texttt{pausasDeCartucho}: fija / lee
		las pausas por plataforma (\texttt{simulador}, \texttt{externo}).
		\item \texttt{fijarAceleracion} / \texttt{aceleracionDeCartucho}:
		fija / lee la aceleración (\texttt{cada}, \texttt{resta},
		\texttt{minimo}).
		\item \texttt{fijarRutas} / \texttt{rutasDeCartucho}: fija / lee la
		política de rutas (\texttt{aumentarCada}).
		\item \texttt{encolar}: encola una pulsación en el buffer FIFO con
		tapa (8); lo descarta y devuelve false si está lleno.
		\item \texttt{colaPulsaciones}: el buffer de pulsaciones pendientes
		(lo vacía el motor al empezar el tick).
	\end{itemize}
	
	\subsection*{baseExtractor.js — extracción de la imagen maestra}
	
	\begin{itemize}
		\item \texttt{verificarBN}: bloqueante si \texttt{img\_s} no es
		blanco y negro puros; lista los colores intrusos.
		\item \texttt{leerImagen}: vuelca una imagen a canvas y devuelve
		sus píxeles (\texttt{getImageData}).
		\item \texttt{extraerBN}: máscara negra + componentes conexas;
		devuelve las componentes elementales, el índice píxel→componente y
		el ruido descartado.
		\item \texttt{componentes}: búsqueda en anchura con pila explícita y
		conectividad-8 sobre la máscara; calcula la caja de cada componente.
		\item \texttt{reagrupar}: aplica las uniones (union-find) y produce
		grupos ordenados de forma determinista y sprites únicos
		deduplicados por huella píxel-exacta.
		\item \texttt{crearUnionFind}: estructura de conjuntos disjuntos
		para las fusiones.
		\item \texttt{buscar}: raíz de una componente en el union-find
		(con compresión de caminos).
		\item \texttt{unir}: une dos componentes en el union-find.
		\item \texttt{hex}: color RGB a notación hexadecimal (para los
		avisos del verificador).
	\end{itemize}
	
	\subsection*{baseGraficos.js — dibujo del mapa de autoría}
	
	\begin{itemize}
		\item \texttt{fijarFondo} / \texttt{fijarSegmentos}: estado interno
		del dibujante: la imagen de fondo y la lista de segmentos residente.
		\item \texttt{pixelesDe}: los píxeles de una imagen como máscara
		B/N.
		\item \texttt{escalaDe}: escala por eje del destino (aparato o
		mapa).
		\item \texttt{DibujarFondo}: pinta la imagen de fondo estirada sobre
		el destino.
		\item \texttt{DibujarSegmento}: pinta un segmento sobre el destino
		(XOR en modo simulador; color de tipo en modo autoría).
	\end{itemize}
	
	\subsection*{baseDescomponer.js — la autoría: mapa, panel, verificación y GWI}
	
	\textbf{Ciclo de vida:}
	\begin{itemize}
		\item \texttt{inicializar}: registra los listeners del mapa y de los
		botones del editor.
		\item \texttt{arrancar}: arranque de la autoría según el tipo:
		GWB $\rightarrow$ descomposición; GWI $\rightarrow$ carga directa de
		listas y edición.
		\item \texttt{reset}: limpieza total de la autoría (selección,
		marcas, historial, ficticios).
		\item \texttt{reconstruir}: mezcla extracción + ficticios +
		historial y produce los grupos finales (hereda props, calcula
		\texttt{oculto}, fusiona por absorción).
		\item \texttt{refrescar} / \texttt{refrescarInfo}: repinta mapa,
		aparato, panel y lista auxiliar.
	\end{itemize}
	
	\textbf{Interacción con el mapa:}
	\begin{itemize}
		\item \texttt{alPulsarMapa}: clic del mapa: aplica la elección
		pendiente, selecciona la pieza o inicia el recuadro de selección.
		\item \texttt{alMoverMapa}: arrastra la banda elástica del recuadro
		de selección múltiple.
		\item \texttt{alSoltarMapa}: cierra el recuadro (marca los grupos
		cortados) o confirma la selección simple.
		\item \texttt{puntoDe}: coordenadas de un evento sobre el canvas.
		\item \texttt{grupoDelPixel}: grupo en las coordenadas del mapa
		(vía el índice píxel$\rightarrow$componente del extractor).
		\item \texttt{grupoDeComp}: grupo actual que contiene una componente.
		\item \texttt{piezaDe}: sube toda la jerarquía de la pieza clicada
		(barra $\rightarrow$ dígito $\rightarrow$ contador).
	\end{itemize}
	
	\textbf{Panel de propiedades:}
	\begin{itemize}
		\item \texttt{mostrarPanel}: pinta el panel según el tipo y la
		selección (incluida la ficha del cartucho).
		\item \texttt{camposDe}: los campos y enlaces que el panel ofrece a
		cada tipo de segmento.
		\item \texttt{aplicarEleccion}: aplica el enlace o unión elegida con
		las validaciones de destino.
		\item \texttt{moverNumero}: Subir/Bajar el número de presentación de
		un segmento (materializa \texttt{orden}, acción atómica).
		\item \texttt{numeroDe}: orden de presentación de los grupos
		(\texttt{orden} o índice).
	\end{itemize}
	
	\textbf{Agrupaciones (unión, dígito, contador):}
	\begin{itemize}
		\item \texttt{unirMarcados}: une por absorción los grupos del
		recuadro (delega en \texttt{hacerDigito} si es un «8»).
		\item \texttt{unirGrupos}: unión por absorción: crea la pieza nueva
		(Fnn), hereda props y re-apunta los enlaces externos.
		\item \texttt{hacerDigito}: clasifica y agrupa las 7 barras de un
		«8» en un dígito (\texttt{barra\_a}\ldots\texttt{barra\_g}).
		\item \texttt{clasificarBarras}: asigna las letras a--g a las barras
		por su geometría (horizontales por altura, verticales por lado).
		\item \texttt{hacerContador}: agrupa de 1 a 4 dígitos en un
		contador (\texttt{digito\_1} = unidades, obligatorio).
		\item \texttt{separarAgrupado}: disuelve una pieza (baja el
		ficticio y anula los enlaces externos que le apuntaban).
		\item \texttt{crearFicticio}: alta de una pieza ficticia (unión,
		dígito o contador) con su número de pieza.
		\item \texttt{barrasDe} / \texttt{digitoDeBarra} /
		\texttt{contadorDeDigito}: navegación de la jerarquía de los
		contadores.
		\item \texttt{esBarraDeDigito} / \texttt{esDigitoDeContador}:
		pertenencia a una agrupación (viajan con \texttt{visible:false}).
		\item \texttt{esAbsorbido} / \texttt{ficticioAbsorbente}: grupo
		absorbido por una unión / la pieza que lo absorbió.
		\item \texttt{fusionarGrupos} / \texttt{imagenFusionada} /
		\texttt{firmaPieza} / \texttt{pixelesDeImagen}: composición OR de
		los píxeles de las piezas fusionadas y su imagen resultante.
		\item \texttt{aportantes}: los grupos que aportan el primer valor no
		nulo de cada prop a la pieza fusionada.
	\end{itemize}
	
	\textbf{Historial (deshacer/rehacer):}
	\begin{itemize}
		\item \texttt{registrarCambio} / \texttt{registrarCambios}: añaden
		entradas al historial (asignan \texttt{n} y \texttt{relacionados};
		vacían la pila de rehacer).
		\item \texttt{deshacer}: revierte el último bloque atómico aplicando
		los \texttt{antes} en orden inverso (fallback por replay en GWI
		antiguos).
		\item \texttt{rehacer}: reaplica el bloque desde la pila de
		rehacibles.
		\item \texttt{invertirEntrada}: aplica el \texttt{antes} de una
		entrada (o la opuesta estructural de unir/separar).
		\item \texttt{quitarParUnion}: retira un par de unión (inversa del
		unir).
		\item \texttt{reproducirCambios}: replay de las entradas del
		historial (reconvierte los \texttt{Fnn} a raíz negativa).
	\end{itemize}
	
	\textbf{Verificación y guardado:}
	\begin{itemize}
		\item \texttt{verificarYGuardar}: la verificación completa
		(bloqueantes y avisos) y el empaquetado del GWI.
		\item \texttt{verificarReferencias}: enlaces colgantes y campos
		obligatorios al cargar un GWI.
		\item \texttt{guardarGWI}: escribe el GWI (campos al nivel
		superior, piezas Fnn, \texttt{visible:false} en lo agrupado).
		\item \texttt{guardarGWIEnDisco}: empaqueta el ZIP con JSZip
		(sprites en \texttt{segmentos/}, \texttt{img\_f}/\texttt{img\_s})
		y fuerza la descarga.
		\item \texttt{formatearGWI}: serializa el GWI compacto (una línea
		por sprite, botón y segmento).
	\end{itemize}
	
	\textbf{Datos del cartucho y del puente:}
	\begin{itemize}
		\item \texttt{cargarAutoria}: carga y reproduce el historial
		\texttt{autoria.cambios} del GWI.
		\item \texttt{cargarListas} / \texttt{segmentosEnMemoria}: vuelcan
		las listas al puente \texttt{comunListas}.
		\item \texttt{asegurarCartucho}: crea la ficha del cartucho (tipo 7)
		si el GWI no la trae.
		\item \texttt{cartuchoDe} / \texttt{esCartucho}: acceso a la ficha
		del cartucho / identificación del tipo 7.
		\item \texttt{pausasDelCartucho} / \texttt{aceleracionDelCartucho}
		/ \texttt{rutasDelCartucho} / \texttt{secuenciaDelCartucho} /
		\texttt{finDelCartucho}: lectura de las propiedades del cartucho
		para la ficha del editor.
		\item \texttt{sincronizarDirector} /
		\texttt{sincronizarSecuencia}: vuelcan en vivo ficha y Secuencia al
		puente en cada reconstrucción (editar = visible para el motor).
		\item \texttt{propPorDefecto} / \texttt{propsDe}: props por defecto
		de una componente / las persistidas de una raíz.
	\end{itemize}
	
	\textbf{Presentación:}
	\begin{itemize}
		\item \texttt{pintarMapa}: el mapa de objetos (cajas, números,
		colores de tipo y de relación, recuadro en curso).
		\item \texttt{listaMotor}: la lista \texttt{\{x, y, imagen, color\}}
		que pintan el mapa y la pantalla de autoría.
		\item \texttt{cajaPresentable}: caja del grupo con el offset aplicado
		(posición presentada, no original).
		\item \texttt{pintarPantallaAutoria}: la pantalla del aparato según
		la fase real de autoría (solo fondo o composición completa).
		\item \texttt{pintarGWI}: composición fondo + segmentos del GWI
		cargado.
		\item \texttt{pintarFondoAparato} / \texttt{pintarSegmentoAparato}:
		fondo estirado / segmento (negro o color de tipo) en la pantalla del
		aparato.
		\item \texttt{fijarEscalaAparato}: escala por eje de la pantalla del
		aparato (\texttt{canvas}/\texttt{img\_s}).
		\item \texttt{colorRGB}: notación hexadecimal a componentes RGB.
		\item \texttt{relacionadosDe}: cadena saliente transitiva del
		seleccionado (azul).
		\item \texttt{entrantesDe}: enlaces entrantes directos al
		seleccionado (magenta).
		\item \texttt{numRuta} / \texttt{numContador}: número de ruta de un
		inicio / de contador de un segmento vida o muerte.
		\item \texttt{segId}: identificador legible de un grupo
		(\texttt{seg\_NNN} o \texttt{Fnn} en los ficticios).
		\item \texttt{resolverEnlace}: índice de componente a su
		\texttt{segId} legible (para el GWI y el panel).
		\item \texttt{listar}: la lista auxiliar de lo numerable (rutas,
		contadores, primer jugador).
		\item \texttt{segmentos}: los segmentos finales (sin fantasmas
		tipo 99).
	\end{itemize}
	
	\subsection*{baseDirector.js — orquestación de máquina y cartucho}
	
	\begin{itemize}
		\item \texttt{inicializar}: registra los listeners globales
		(botones, ficheros, teclado).
		\item \texttt{alSeleccionarSetup} / \texttt{alSeleccionarCartucho}
		/ \texttt{alSeleccionarGWIZip}: los selectores de fichero; validan
		y lanzan el flujo.
		\item \texttt{validarSetup} / \texttt{validarCartucho}: validación
		bloqueante del setup (estilo, pantalla, botones) y del cartucho
		(tipo y versión obligatorios, estilo compatible).
		\item \texttt{aplicarSetup}: construye la máquina del estilo y
		muestra el canvas en su hueco.
		\item \texttt{aplicarBotonesCartucho}: renombra los botones usados
		con su \texttt{id} y elimina del DOM los no usados
		(setup $\geq$ cartucho).
		\item \texttt{construirBotonera}: esqueleto del estilo con huecos
		anclados en rejilla (sin re-flujos).
		\item \texttt{crearBotonFisico}: nodo DOM de un botón físico con su
		número, alias y tecla.
		\item \texttt{procesarCartucho}: el flujo según el tipo de
		cartucho: GWB $\rightarrow$ descomposición; GWI $\rightarrow$ carga
		directa; GWC $\rightarrow$ sin procesar (futuro ESP32).
		\item \texttt{elegirJSONCartucho}: elige el JSON de la carpeta por
		prioridad GWC $>$ GWI $>$ GWB.
		\item \texttt{descomponerImagen}: \texttt{verificarBN} +
		\texttt{extraerBN} + arranque de la autoría.
		\item \texttt{construirMapaFicheros}: la lista de ficheros a un mapa
		nombre $\rightarrow$ File.
		\item \texttt{leerJSON}: lectura y \texttt{parse} de un JSON con
		mensajes de error legibles.
		\item \texttt{cargarImagen} / \texttt{cargarImagenAux} /
		\texttt{imagenBlob} / \texttt{crearURL}: carga de imágenes desde
		fichero o blob (con object URLs revocables).
		\item \texttt{validarTamanoImgS}: valida \texttt{img\_s} contra
		\texttt{img\_f} (aviso si difieren).
		\item \texttt{pintarFondoCartucho} / \texttt{pintarFondoNeutro}:
		\texttt{img\_f} estirada en la pantalla / pantalla sin cartucho.
		\item \texttt{mostrarSprites}: tira de miniaturas con su id bajo la
		máquina.
		\item \texttt{mostrarListaFicheros}: lista de ficheros leídos con
		tamaños en el panel inferior.
		\item \texttt{resetInfoCartucho} / \texttt{limpiarCartucho}: limpian
		la línea de cartucho del panel y revocan las URLs al recargar.
		\item \texttt{avisar} / \texttt{avisarError} / \texttt{avisarEtiquetas}:
		la barra de estado (última operación, errores en rojo).
		\item \texttt{normalizarNumeroBoton}: acepta «A»/«6»/números sin
		comillas y devuelve el número físico 0--8.
		\item \texttt{alPulsarTecla} / \texttt{alSoltarTecla}: el teclado de
		la máquina: resaltado del botón y reenvío al \texttt{juegoDirector}
		(con el juego en marcha).
		\item \texttt{buscarBotonPorTecla}: tecla $\rightarrow$ botón físico.
		\item \texttt{alAtajoTeclado}: los atajos de edición
		(Ctrl+Z/Ctrl+Y deshacer/rehacer).
		\item \texttt{jugar}: el puente del botón Start hacia el
		\texttt{juegoDirector}.
		\item \texttt{comprobarSalud}: diagnóstico: ficheros cargados,
		globales esperados y avisos.
		\item \texttt{verificarYGuardar}: delega en
		\texttt{Descomponer.verificarYGuardar} el botón de guardado.
	\end{itemize}
	
	\subsection*{juegoDirector.js — el reloj del simulador}
	
	\begin{itemize}
		\item \texttt{arrancar}: relanza la partida: reinicia semilla, cola,
		contadores de rutas y llama a \texttt{iniciar()} del motor; arranca
		el reloj.
		\item \texttt{programar}: programa el siguiente tick con la
		\textbf{pausa vigente} (\texttt{setTimeout} encadenado, nunca
		\texttt{setInterval}).
		\item \texttt{detener} / \texttt{enMarcha}: detiene el reloj /
		pregunta si el juego está en marcha.
		\item \texttt{tick}: un tick de reloj: pasa al motor y actualiza la
		barra de diagnóstico (ticks, vidas, muertes, Secuencia, pausa).
		\item \texttt{pulsar}: acepta una pulsación solo con el juego en
		marcha y la encola (el director no interpreta entradas).
		\item \texttt{pausaVigente}: la pausa actual aplicando la
		aceleración (nunca por debajo del mínimo).
		\item \texttt{aleatorio}: el generador LCG de 16 bits: única puerta
		del azar, rango $[0, max)$.
		\item \texttt{reiniciarAleatorios}: reinicia el estado del LCG a la
		semilla (misma semilla, misma partida).
		\item \texttt{rutaLanzada}: cuenta el total de rutas lanzadas
		(monótono, contador de progresión).
		\item \texttt{maxRutasSimultaneas}: máximo de rutas simultáneas (+1
		en cada múltiplo de \texttt{aumentarCada}, tope el número de rutas).
	\end{itemize}
	
	\subsection*{juegoMotorEstados.js — el motor de estados del juego}
	
	\textbf{Ciclo de vida:}
	\begin{itemize}
		\item \texttt{iniciar}: restaura la partida como el cartucho la
		dibujó: posiciones originales, contadores, familias de iconos y
		rutas.
		\item \texttt{lista}: pide la lista de segmentos al puente en cada
		uso (nunca guarda una referencia muerta).
		\item \texttt{tick}: el orden sagrado: (0) entrada, (1) apagar
		banderas \texttt{cambiar}, (2) barrido de rutas y desenlaces; tras
		él evalúa el fin de partida.
		\item \texttt{terminada} / \texttt{congelada}: estado del juego
		partida terminada / congelada por cadena de desenlace.
	\end{itemize}
	
	\textbf{Entrada y jugador:}
	\begin{itemize}
		\item \texttt{pulsar}: recibe el número físico del botón (nunca un
		evento del DOM).
		\item \texttt{aplicarEntrada}: vacía el buffer del tick y mueve al
		jugador por los enlaces de dirección.
		\item \texttt{combinar}: reduce el buffer a ejes cuando la regla
		está activa (un solo paso por tick, diagonales permitidas).
		\item \texttt{moverJugador}: mueve por el enlace de dirección
		(escribir en la lista: posición y bandera de cambio).
		\item \texttt{jugadorActual} / \texttt{buscarJugador}: el índice del
		segmento jugador / su búsqueda (marca del diseñador primero,
		salvavidas después).
		\item \texttt{indiceDe}: localiza un segmento por su id
		(\texttt{seg\_NNN}).
	\end{itemize}
	
	\textbf{Rutas y colisiones:}
	\begin{itemize}
		\item \texttt{esDeRuta} / \texttt{esCadena}: pertenencia a rutas
		(tipos juego/inicio) / a cadenas de desenlace (todo menos jugador).
		\item \texttt{lanzarRuta}: lanza si hay hueco según el máximo de
		simultáneas (elige ruta con el PRNG y su lista de repeticiones).
		\item \texttt{rutasVivas} / \texttt{estadoRutas} / \texttt{muereRuta}:
		el número de rutas vivas / el estado completo de rutas / la baja de
		una ruta (cinco puntos de muerte).
		\item \texttt{decidirColision}: el corazón: lee la bandera
		\texttt{visible} del jugador señalado y sigue \texttt{salvado} o
		\texttt{muerte} (el contador indicado sube una sola vez, en la
		decisión).
	\end{itemize}
	
	\textbf{Contadores y familias de iconos:}
	\begin{itemize}
		\item \texttt{prepararContadores}: localiza los contadores de
		dígitos del cartucho (qué segmentos los componen).
		\item \texttt{ajustarContador}: fija el valor de uno de los cuatro
		contadores internos.
		\item \texttt{numeroDeContador}: el índice de contador (0--3) que
		declara el desenlace de una colisión.
		\item \texttt{familiaIconos} / \texttt{esIconoFamilia} /
		\texttt{ordenFamilia}: la familia de iconos vida/muerte / su
		identificación / su orden por la cadena de ruta.
		\item \texttt{sincronizarIconos} / \texttt{setVis}: enciende o apaga
		cada icono de la familia según el contador correspondiente.
		\item \texttt{contarAterrizaje}: acumula en el icono que le toca por
		el orden de la familia.
		\item \texttt{filaDe}: las barras de un dígito por su letra (a--g).
		\item \texttt{presentarContadores}: enciende en la pantalla los
		dígitos de cada contador (tabla de siete segmentos, rellenos).
		\item \texttt{contadores}: los contadores internos para la barra de
		diagnóstico.
	\end{itemize}
	
	\textbf{Animaciones y desenlaces:}
	\begin{itemize}
		\item \texttt{avanzarSecuencia}: la Secuencia de fondo: un segmento
		por tick y el ciclo $n + m$.
		\item \texttt{secuenciaInfo}: el estado de la Secuencia para la
		barra de diagnóstico.
		\item \texttt{avanzarAdicionales}: el bucle de los elementos sin
		interacción (el helicóptero, las aspas).
		\item \texttt{sincronizarAcompanantes}: enciende o apaga el
		secundario «Mostrado junto a» siempre con su principal.
		\item \texttt{comprobarFin}: evalúa el mecanismo de fin elegido
		(contador a cero, iconos de vida o de muerte completos).
		\item \texttt{debug} / \texttt{pruebaAleatorios}: diagnóstico bajo
		GW\_DEBUG / la función de prueba del generador (10 números, dos
		pasadas idénticas).
	\end{itemize}
	
	\subsection*{juegoMotorGrafico.js — el motor gráfico}
	
	\begin{itemize}
		\item \texttt{inicializar}: pinta solo el fondo; avisa si no hay
		cartucho.
		\item \texttt{tick}: el repintado completo de la escena: fondo + todos
		los segmentos visibles (lectura pura de la lista).
		\item \texttt{pintarFondo}: blanco + \texttt{img\_f} estirado (del
		puente).
		\item \texttt{pintarSegmento}: un segmento en negro puro, escalado
		por vecino más cercano y por eje, con la caja recortada al canvas.
		\item \texttt{pixelesDe}: los píxeles del sprite como máscara.
		\item \texttt{debug}: diagnóstico bajo GW\_DEBUG.
	\end{itemize}
	
	\clearpage
	
\end{document}